% ===================================================================
%  图表第 6 号 — 屋子里面。— 家具。— 饮食。— 照明。
%  Traduction 2026 d'apres texto/io, controlee sur texto/fr, texto/en
%  et texto/es. Consigne : voir texto/zh/10-tubiao-01.tex.
%
%  Deux notes, toutes deux marquees « (*) », et deux appels : celui de
%  « babo » au bloc c1-12-1, celui de « lambrequino » au bloc c2-01-1.
%  L'ordre les departage, comme dans les autres colonnes.
%
%  LA FEMME DE CHAMBRE. L'ido du bloc c1-06-1 porte « (6) » la ou le
%  francais porte « (16) », et c'est le francais qui a raison. Le
%  chinois suit le francais, comme l'anglais, l'espagnol et le russe.
% ===================================================================

%%K t06-apar-1 apar
\VUsaut{6.0mm}
\VUtitre{15.0pt}{60}{图表第 6 号}
\VUsaut{1.4mm}
\VUtitre{11.4pt}{0}{屋子里面、家具、饮食、}
\VUsaut{0.4mm}
\VUtitre{11.4pt}{0}{照明。}
\VUsaut{2.2mm}

%%K t06-apar-2 apar
房子盖好了。约翰的叔叔搬进了他的新居，那布置我们已经知道。现在来看看他是
怎样陈设的。我们先去看：

%%K t06-tit-1 sub
\VUpk{10.2pt}{40}{卧房。}
\VUsaut{3.0mm}

%%K t06-c1-01-1 p
1. --- 我们来得不巧，因为女仆正在收拾\VUgras{房间}。

%%K t06-c1-02-1 p
2. --- \VUgras{洗脸架} \textsuperscript{(1)} 上东一样西一样地放着洗漱梳妆用的
东西。那是\VUgras{脸盆} \textsuperscript{(2)} 和\VUgras{水壶} \textsuperscript{(3)}、
\VUgras{头刷} \textsuperscript{(4)}、\VUgras{梳子} \textsuperscript{(5)}、
\VUgras{肥皂} \textsuperscript{(6)} 和\VUgras{海绵} \textsuperscript{(7)}，最后还有
一小\VUgras{瓶} \textsuperscript{(8)} 漱口水。

%%K t06-c1-03-1 p
3. --- 地上，洗脸架前面，是盛脏水的\VUgras{桶} \textsuperscript{(9)}。

%%K t06-c1-04-1 p
4. --- \VUgras{外间} \textsuperscript{(10)} 里看得见几个\VUgras{衣帽钩}
\textsuperscript{(13)} 和两把\VUgras{伞} \textsuperscript{(11)}，插在\VUgras{伞架}
\textsuperscript{(12)} 里。

%%K t06-c1-05-1 p
5. --- 门那边立着\VUgras{带镜衣橱} \textsuperscript{(14)}，里面放着\VUgras{衣物}
\textsuperscript{(15)}。

%%K t06-c1-06-1 p
6. --- 趁\VUgras{女仆} \textsuperscript{(16)} 整\VUgras{床} \textsuperscript{(17)}
的工夫，我们来看看床的各个部分。\VUgras{床架} \textsuperscript{(20)} 上是一张
好\VUgras{弹簧褥} \textsuperscript{(21)}，朱斯蒂娜等会儿要在上面铺一床毛
\VUgras{褥子} \textsuperscript{(22)}。然后她从椅子上取来两条\VUgras{床单}
\textsuperscript{(23)} 和一床\VUgras{毯子} \textsuperscript{(24)}。再把\VUgras{长枕}
\textsuperscript{(25)} 和\VUgras{枕头} \textsuperscript{(26)} 放到床头，它们这会儿
同\VUgras{鸭绒被} \textsuperscript{(27)} 一起放在\VUgras{床前小毯} \textsuperscript{(32)}
上。她用鸡毛掸子掸\VUgras{床帷} \textsuperscript{(19)} 和\VUgras{床顶}
\textsuperscript{(18)}，活儿就做了一半。

%%K t06-c1-07-1 p
7. --- 然后她还得把\VUgras{床头柜} \textsuperscript{(28)} 略微收拾收拾，给
\VUgras{闹钟} \textsuperscript{(29)} 上弦，把\VUgras{夜灯} \textsuperscript{(30)}
预备好，再把\VUgras{蜡烛} \textsuperscript{(31)} 拿走，好擦烛台。

%%K t06-tit-2 sub
\VUsaut{5.0mm}
\VUpk{10.2pt}{40}{针线篮和壁炉用具。}
\VUsaut{3.0mm}

%%K t06-c1-08-1 p
8. --- \VUgras{针线篮} \textsuperscript{(34)} 放在\VUgras{小凳} \textsuperscript{(33)}
旁边，也很该收拾了。她得把\VUgras{绣活} \textsuperscript{(35)} 叠起来，把
\VUgras{针线桌} \textsuperscript{(38)} 收拾好——\VUgras{顶针}、\VUgras{线}
\textsuperscript{(36)}、\VUgras{针} \textsuperscript{(37)} 和\VUgras{剪刀}
\textsuperscript{(39)} 都收进去——再用\VUgras{钥匙} \textsuperscript{(40)} 把桌盖
锁上。

%%K t06-c1-09-1 p
9. --- 屋子当中的\VUgras{圆桌} \textsuperscript{(41)} 上，朱斯蒂娜要给
\VUgras{玻璃瓶} \textsuperscript{(42)} 换水。她要往\VUgras{糖罐} \textsuperscript{(43)}
里放\VUgras{糖}，擦\VUgras{糖夹} \textsuperscript{(44)}，把\VUgras{衣刷}
\textsuperscript{(46)} 拿走。

%%K t06-c1-10-1 p
10. --- 屋子右边给她的活儿少些，因为\VUgras{卧榻} \textsuperscript{(47)} 和它的
\VUgras{垫子} \textsuperscript{(48)} 都在原处，天又不冷，\VUgras{壁炉}
\textsuperscript{(49)} 那一套东西也没人动过。\VUgras{火铲} \textsuperscript{(50)}、
\VUgras{火钳} \textsuperscript{(51)}、\VUgras{柴架} \textsuperscript{(55)} 和
\VUgras{挡灰栅} \textsuperscript{(56)} 都很干净；\VUgras{炉屏} \textsuperscript{(54)}
没有挪动，\VUgras{柴箱} \textsuperscript{(52)} 里\VUgras{劈柴} \textsuperscript{(53)}
装得满满的。

%%K t06-noto-1 noto
\VUnotes{143.2mm}{%
(*) Babo = 小孩子；英语 \textit{baby}，法语 \textit{bébé}，意大利语
\textit{bambino}。%
}

%%K t06-noto-2 noto
\VUnotes{143.2mm}{%
(*) Lambrequino = 法语 \textit{lambrequin}，西班牙语
\textit{lambrequines}，葡萄牙语 \textit{lambrequins}，连德语和英语也说
\textit{lambrequins} ——所以德、英、法、葡、西五种语言都一样。%
}

%%K t06-c1-11-1 p
11. --- 她还得掸\VUgras{座钟} \textsuperscript{(57)}、\VUgras{烛台}
\textsuperscript{(58)} 和\VUgras{小托盘} \textsuperscript{(59)}，再擦\VUgras{镜子}
\textsuperscript{(60)}。

%%K t06-c1-12-1 p
12. --- 至于娃娃的\VUgras{摇篮} \textsuperscript{(61)}（就是那个 babo (*)），
已经预备好了。

%%K t06-tit-3 sub
\VUsaut{5.0mm}
\VUpk{10.2pt}{40}{客厅。}
\VUsaut{3.0mm}

%%K t06-c2-01-1 p
1. --- 现在是客厅了。这是一间很漂亮的\VUgras{屋子}。右角上的壁炉摆着一尊
精巧的\VUgras{小像} \textsuperscript{(1)} 和两只好看的\VUgras{花瓶}
\textsuperscript{(3)}，还垂着一幅丝绒\VUgras{壁炉帷} \textsuperscript{(2)} (*)。

%%K t06-c2-02-1 p
2. --- 为了不让炉膛里的火星把地毯点着，炉前放了一面铜\VUgras{火挡}
\textsuperscript{(4)}。

%%K t06-c2-03-1 p
3. --- 旁边敞着一张精致的女用\VUgras{写字台} \textsuperscript{(5)}。
\VUgras{女主人} \textsuperscript{(23)} 就在这里写信。

%%K t06-c2-04-1 p
4. --- 窗上挂着\VUgras{纱帘} \textsuperscript{(6)}，用\VUgras{束带}
\textsuperscript{(8)} 拢着，束带挂在\VUgras{挂钩} \textsuperscript{(7)} 上；还有
厚布窗帘，上面加着一道\VUgras{帘楣} \textsuperscript{(9)}。

%%K t06-c2-05-1 p
5. --- 窗洞里一只\VUgras{花架} \textsuperscript{(11)} 托着一盆绿\VUgras{植物}
\textsuperscript{(10)}，那是一株很好的棕榈，叶子几乎够到了\VUgras{煤油灯}
\textsuperscript{(12)}。

%%K t06-c2-06-1 p
6. --- \VUgras{灯罩} \textsuperscript{(13)} 是玛丽做的，就是那位可爱的
\VUgras{小姐} \textsuperscript{(14)}，坐在\VUgras{钢琴} \textsuperscript{(15)} 旁。她笔直地坐在
\VUgras{琴凳} \textsuperscript{(16)} 上练一支曲子。她很有本事，也很整齐，看她的
\VUgras{琴谱柜} \textsuperscript{(17)} 就知道：里面井井有条。

%%K t06-tit-4 sub
\VUsaut{5.0mm}
\VUpk{10.2pt}{40}{淘气鬼。}
\VUsaut{3.0mm}

%%K t06-c2-07-1 p
7. --- 她弟弟保罗在找一本\VUgras{书} \textsuperscript{(19)}，在\VUgras{书橱}
\textsuperscript{(18)} 里翻。这\VUgras{孩子} \textsuperscript{(20)} 坐不住，很叫人受不了。
他吊在书橱上去够那本放得太高的书，眼看要把橱顶上的\VUgras{半身像}
\textsuperscript{(21)} 弄下来。

%%K t06-c2-08-1 p
8. --- 他敢干这蠢事，是趁着母亲正忙着跟一位女友说话，那位\VUgras{太太}
\textsuperscript{(22)} 来同她住几天。母亲坐在\VUgras{长沙发} \textsuperscript{(24)} 上，就在
\VUgras{扶手椅} \textsuperscript{(25)} 旁边，她的女友坐在\VUgras{椅子}
\textsuperscript{(27)} 上，我们只看得见\VUgras{椅背} \textsuperscript{(26)}。

%%K t06-c2-09-1 p
9. --- 墙上挂着的大\VUgras{镜框} \textsuperscript{(30)} 里是一位亲戚的
\VUgras{相} \textsuperscript{(29)}。桌上摊着的\VUgras{相册} \textsuperscript{(32)} 里
收着家里其余人的\VUgras{照片} \textsuperscript{(33)}。孩子们刚才还在翻它，因为
\VUgras{椅子} \textsuperscript{(31)} 还围在桌子四周。

%%K t06-c2-10-1 p
10. --- 那只瓷\VUgras{中国花瓶} \textsuperscript{(34)} 就在\VUgras{裁纸刀}
\textsuperscript{(35)} 旁边，是玛丽命名日的礼物，占着屋角的那架\VUgras{屏风}
\textsuperscript{(36)} 是她叔叔从中国带回来的。

%%K t06-c2-11-1 p
11. --- 漂亮的\VUgras{画} \textsuperscript{(37)} 挂在\VUgras{墙}
\textsuperscript{(42)} 上，使它明快，天花板上的\VUgras{吊灯} \textsuperscript{(38)} 给它
照明，那些\VUgras{电灯泡} \textsuperscript{(39)} 一到晚上就亮得叫人高兴，这间
客厅看着很舒服。地上铺的软\VUgras{地毯} \textsuperscript{(40)} 和那么方便的
\VUgras{电铃} \textsuperscript{(41)}，使它十分舒适。

%%K t06-tit-5 sub
\VUsaut{3.6mm}
\VUpk{10.2pt}{40}{饭厅。}
\VUsaut{3.0mm}

%%K t06-c3-01-1 p
1. --- 开饭的铃响了。除了玛丽，全家都围在桌旁。饭厅里有一个很好的陶
\VUgras{炉子} \textsuperscript{(1)}，带一根细\VUgras{烟囱} \textsuperscript{(2)}，
烘得很暖和。

%%K t06-c3-02-1 p
2. --- 这间屋子家具不多。靠墙，\VUgras{小凳} \textsuperscript{(4)} 上方，挂着
一个好看的小\VUgras{洗手器} \textsuperscript{(3)}。旁边是\VUgras{餐具柜}
\textsuperscript{(5)}，上面放着冷菜、甜点和一时用不着的各样餐具。

%%K t06-c3-03-1 p
3. --- 所以上一层看得见一只\VUgras{烤鸡} \textsuperscript{(7)}、一条\VUgras{鱼}
\textsuperscript{(8)} 和一\VUgras{盘} \textsuperscript{(10)} \VUgras{水果}
\textsuperscript{(9)}。下面是\VUgras{奶酪} \textsuperscript{(11)}、\VUgras{面包}
\textsuperscript{(12)} 和\VUgras{调味架} \textsuperscript{(13)}，架上有拌沙拉要用的
一切：油、醋、\VUgras{盐} \textsuperscript{(14)} 和\VUgras{胡椒} \textsuperscript{(15)}。
最后，最下一层放着一个\VUgras{蛋糕} \textsuperscript{(17)}，旁边是\VUgras{芥末罐}
\textsuperscript{(16)}。

%%K t06-c3-04-1 p
4. --- 饭厅的钟叫\VUgras{挂钟} \textsuperscript{(20)}，样子很别致。它嵌在一个
木框里，框里还有\VUgras{气压表} \textsuperscript{(19)} 和\VUgras{温度表}
\textsuperscript{(18)}。

%%K t06-tit-6 sub
\VUsaut{4.6mm}
\VUpk{10.2pt}{40}{上菜。}
\VUsaut{3.0mm}

%%K t06-c3-05-1 p
5. --- 屋子四周墙上都镶着打蜡橡木的\VUgras{护壁板} \textsuperscript{(21)}。一幅
布\VUgras{门帘} \textsuperscript{(22)} 把通往凉廊的门遮得挺严。饭后全家就到那里
去喝咖啡。\VUgras{杯子} \textsuperscript{(24)} 和\VUgras{碟子} \textsuperscript{(25)}
已经摆在\VUgras{瓷器柜} \textsuperscript{(23)} 上了。

%%K t06-c3-06-1 p
6. --- 桌子上方天花板上装着一盏\VUgras{吊灯} \textsuperscript{(26)}；晚上它很亮
的光洒满整个饭厅。

%%K t06-c3-07-1 p
7. --- 现在来看\VUgras{餐桌} \textsuperscript{(27)} 本身。全家进来时桌子已经摆
好了。\VUgras{桌布} \textsuperscript{(28)} 上，每个人的位子前放着一只\VUgras{盘子}
\textsuperscript{(33)}、一条\VUgras{餐巾} \textsuperscript{(29)}、一把\VUgras{汤匙}
\textsuperscript{(34)}、一把\VUgras{叉} \textsuperscript{(35)}、一把\VUgras{刀}
\textsuperscript{(36)} 和一只\VUgras{杯} \textsuperscript{(32)}。还有盛酒的
\VUgras{瓶} \textsuperscript{(30)} 和盛水的\VUgras{玻璃瓶} \textsuperscript{(31)}。

%%K t06-c3-08-1 p
8. --- \VUgras{仆人} \textsuperscript{(39)} 先端来了\VUgras{汤盆} \textsuperscript{(37)}，
里面的汤还冒着热气。这会儿他在上第一道\VUgras{菜} \textsuperscript{(38)}，是荤的。
然后他要把我们已经说过的那些端上来；最后，上完\VUgras{甜点}，他要趁主人们
到凉廊去了，把桌子收拾干净，把饭厅重新归置好。

%%K t06-c3-08-2 p
\textit{附注。} --- \textit{关于饮食的日常词语这里只举了大略。这份单子
将在「旅馆」（第 13 号）和「市场」（第 15 号）两幅图表里补全。}

%%K t06-tit-7 sub
\VUsaut{4.6mm}
\VUpk{10.2pt}{40}{厨房。}
\VUsaut{3.0mm}

%%K t06-c4-01-1 p
1. --- 我们从半开的门里望进厨房，看见的是一片好看的杂乱。\VUgras{灶}
\textsuperscript{(1)} 前\VUgras{火} \textsuperscript{(3)} 噼啪响着，架着\VUgras{烤架}
\textsuperscript{(5)}。一大块\VUgras{烤肉} \textsuperscript{(7)} 穿在\VUgras{铁扦}
\textsuperscript{(6)} 上，快烤好了。

%%K t06-c4-02-1 p
2. --- 刚用过的\VUgras{风箱} \textsuperscript{(8)} 和\VUgras{煤铲} \textsuperscript{(9)}
连同\VUgras{柴块} \textsuperscript{(10)} 一起丢在地上。

%%K t06-c4-03-1 p
3. --- 炉台上倒整齐：\VUgras{灯笼} \textsuperscript{(11)}、\VUgras{火柴盒架}
\textsuperscript{(13)} 和里面的\VUgras{火柴} \textsuperscript{(12)}、\VUgras{烛台}
\textsuperscript{(14)} 和\VUgras{手烛台} \textsuperscript{(15)} 连同它的
\VUgras{蜡烛} \textsuperscript{(16)}，都在各自的地方。可是那只大\VUgras{煤桶}
\textsuperscript{(18)}，装满了\VUgras{煤} \textsuperscript{(17)}——煤是\VUgras{厨娘}
\textsuperscript{(23)} 刚下地窖取来的——却丢在厨房当中。

%%K t06-c4-04-1 p
4. --- 说实在的，这可怜的女人得赶紧，午饭才来得及。这会儿她在\VUgras{炉灶}
\textsuperscript{(19)} 前搅一锅\VUgras{调味汁} \textsuperscript{(24)}，那可不能烧糊了。

%%K t06-c4-05-1 p
5. --- \VUgras{烤箱} \textsuperscript{(20)} 里烤着一个蛋糕，是她给孩子们做的
点心。炉灶上方挂着\VUgras{长柄勺} \textsuperscript{(21)} 和\VUgras{漏勺}
\textsuperscript{(22)}。

%%K t06-tit-8 sub
\VUsaut{4.0mm}
\VUpk{10.2pt}{40}{厨房用具。}
\VUsaut{3.0mm}

%%K t06-c4-06-1 p
6. --- 挂在屋子里头的那\VUgras{套锅} \textsuperscript{(25)} 很干净。好看的铜
\VUgras{炖锅} \textsuperscript{(26)}、\VUgras{奶壶} \textsuperscript{(27)}、
\VUgras{滤器} \textsuperscript{(28)} 和\VUgras{擦床} \textsuperscript{(29)} 都仔细
擦洗过。

%%K t06-c4-07-1 p
7. --- \VUgras{盐罐} \textsuperscript{(30)} 放在小碗橱上，旁边是\VUgras{茶壶}
\textsuperscript{(31)} 和\VUgras{油坛} \textsuperscript{(32)}。\VUgras{抽屉}
\textsuperscript{(33)} 里大概放着餐具和菜刀。

%%K t06-c4-08-1 p
8. --- \VUgras{扫帚} \textsuperscript{(34)} 和\VUgras{鸡毛掸子} \textsuperscript{(35)}
在角落里。厨娘刚用过\VUgras{剁墩} \textsuperscript{(36)}；她把它撂在窗子旁边。

%%K t06-c4-09-1 p
9. --- 墙上挂着一条\VUgras{手巾} \textsuperscript{(37)}，旁边是\VUgras{漏斗}
\textsuperscript{(38)}。厨房这一头放着\VUgras{烤网} \textsuperscript{(39)}、
\VUgras{剁刀} \textsuperscript{(40)} 和\VUgras{切菜板} \textsuperscript{(41)}。再往上，
剁刀上方的\VUgras{搁板} \textsuperscript{(42)} 上，是\VUgras{罐子} \textsuperscript{(43)}、
\VUgras{炖罐} \textsuperscript{(44)} 连同它的\VUgras{盖} \textsuperscript{(45)}，
还有\VUgras{大锅} \textsuperscript{(46)}。\VUgras{煎锅} \textsuperscript{(47)} 就挂在近旁。

%%K t06-c4-10-1 p
10. --- 这个角落里只有那只铜\VUgras{龙头} \textsuperscript{(48)} 闪着光。\VUgras{水槽} \textsuperscript{(49)}
上放着大
\VUgras{水罐} \textsuperscript{(50)}，它一定很方便，下面有个小柜，里面放着
\VUgras{醋桶} \textsuperscript{(51)} 连同它的\VUgras{小龙头} \textsuperscript{(52)}、
\VUgras{垃圾箱} \textsuperscript{(53)} 和\VUgras{脏碗碟} \textsuperscript{(54)}。

%%K t06-tit-9 sub
\VUsaut{4.0mm}
\VUpk{10.2pt}{40}{厨房里还放着什么。}
\VUsaut{3.0mm}

%%K t06-c4-11-1 p
11. --- 那只大\VUgras{木盆} \textsuperscript{(55)}，是洗\VUgras{蔬菜}
\textsuperscript{(58)} 用的，还有那口生铁\VUgras{锅} \textsuperscript{(56)}，立在
\VUgras{瓷砖地} \textsuperscript{(57)} 上，大概也是那个柜子里的。

%%K t06-c4-12-1 p
12. --- 厨娘的炉子肯定不缺，因为这边桌角上就有一个\VUgras{煤气灶}
\textsuperscript{(59)}，上面一只小锅正烧着水。冒出来的\VUgras{汽} \textsuperscript{(60)}
告诉我们水该开了。这水大概是煮咖啡的，因为桌子那一头就放着\VUgras{咖啡壶}
\textsuperscript{(64)} 和\VUgras{咖啡磨} \textsuperscript{(63)}。

%%K t06-c4-13-1 p
13. --- 煤气灶接在\VUgras{煤气嘴} \textsuperscript{(62)} 上，用的是一根长长的
\VUgras{橡皮管} \textsuperscript{(61)}。

%%K t06-c4-14-1 p
14. --- 来帮厨娘的\VUgras{短工} \textsuperscript{(66)} 正在收拾晚饭的菜。削好
洗净以后，她要把菜放进\VUgras{锅} \textsuperscript{(65)} 里，等着下锅的时候。

%%K t06-tit-10 sub
\VUsaut{4.0mm}
{\centering\fontsize{10.2pt}{10.2pt}\selectfont\textls[40]{\scshape 浴室。} \textit{（洗澡的屋子。）}\par}
\VUsaut{3.0mm}

%%K t06-c5-01-1 p
1. --- 要认全这所房子的主要屋子，只剩最后一处要探一探了：看看浴室的
设备。这用不了多久，因为它不大，我们很快就看得出\VUgras{澡盆}
\textsuperscript{(1)} 配着一个好\VUgras{热水器} \textsuperscript{(2)}，\VUgras{肥皂碟}
\textsuperscript{(3)} 就在洗澡人手边，\VUgras{毛巾架} \textsuperscript{(5)} 一定使
\VUgras{毛巾} \textsuperscript{(4)} 容易晾干。

%%K t06-c5-02-1 p
2. --- 这间屋子墙上贴着\VUgras{釉面砖} \textsuperscript{(6)}，所以装了
\VUgras{淋浴} \textsuperscript{(7)} 也不怕用的时候溅的水把墙弄坏。一只小锌
\VUgras{盆} \textsuperscript{(8)} 用来洗脚，就把陈设配齐了。
\VUsaut{6.0mm}
\VUfilet{22mm}
